\chapter{Zero as the Critical Axis of the Relational Barycenter}
\label{chap:zero-critical-relational-axis}

\section{Purpose}

This chapter compresses the relational-zero programme to its shortest current
proof graph.  It proves the same half-axis zero from three independent
directions: exchange symmetry, Shannon/KL information geometry, and the
projective reciprocal--conjugation defect.  It then states the exact
operator-facing endgame for the Riemann Hypothesis.

The word \emph{barycenter} below means the exchange-symmetric affine midpoint
of a normalized complementary pair.  It is not an arbitrary mass-weighted
center of unequal physical bodies.  Likewise, \emph{relational singular
node} means the unique zero-defect state selected by the declared relation;
it does not mean an analytic pole.

\section{Minimal relational core}

For every admitted relation let
\[
p_e=(\sigma_e,1-\sigma_e),\qquad
J_e(\sigma_e)=1-\sigma_e,
\qquad
x_e=\sigma_e-\frac12.
\]
Define
\[
D_H(\sigma)
=
\ln2-H_2(\sigma)
=
D_{\mathrm{KL}}\!\left(
(\sigma,1-\sigma)
\middle\|
(1/2,1/2)
\right).
\]
For a finite or countable relation family define the non-negative relational
action-defect
\[
\mathfrak S_{\rm rel}
=
\sum_{e}w_e
\left(
\|F_e\|_{G_e}^2
+
\lambda_eD_H(\sigma_e)
\right),
\qquad w_e,\lambda_e>0,
\]
where every \(G_e\) is positive definite.

\begin{warning}
The project quantity \(\mathfrak S_{\rm rel}\) is not the ordinary
Lorentzian action \(S=\int L\,dt\).  Its zero condition must not be
identified with the standard variational condition \(\delta S=0\).
\end{warning}

\section{First side: exchange-symmetric relational center}

\begin{theorem}[Fixed-point center]
For every normalized complementary relation,
\[
J_e(\sigma_e)=\sigma_e
\iff
\sigma_e=\frac12
\iff
x_e=0.
\]
\end{theorem}

\begin{proof}
The fixed-point equation is \(1-\sigma_e=\sigma_e\), hence
\(2\sigma_e=1\).  The centered-coordinate statement is immediate.
\end{proof}

For an affine pair \((a_e,b_e)\), the exchange-symmetric midpoint is
\[
c_e=\frac{a_e+b_e}{2}.
\]
Under affine normalization \(a_e\mapsto0\), \(b_e\mapsto1\), this becomes
\(c_e=1/2\).  Thus the half is the normalized relational barycenter.

\section{Second side: Shannon/KL information zero}

\begin{theorem}[Information center]
For \(0<\sigma<1\),
\[
D_H(\sigma)\ge0
\]
and
\[
D_H(\sigma)=0
\iff
\sigma=\frac12
\iff
H_2(\sigma)=\ln2.
\]
Moreover \(1/2\) is the unique stationary point of \(D_H\).
\end{theorem}

\begin{proof}
Non-negativity and the equality condition follow from the KL-divergence
identity.  Independently,
\[
D_H'(\sigma)=\ln\frac{\sigma}{1-\sigma},
\qquad
D_H''(\sigma)=\frac1\sigma+\frac1{1-\sigma}>0.
\]
Hence the sole stationary point is \(\sigma=1/2\), and it is the strict
global minimum of the defect.
\end{proof}

Near the center, with \(x=\sigma-1/2\),
\[
D_H\!\left(\frac12+x\right)
=
2x^2+\frac43x^4+O(x^6),
\]
so every nonzero transverse displacement carries positive information defect.

\section{Relational Lagrange closure}

\begin{theorem}[Relational Lagrange--Zero Theorem]
\[
\boxed{
\mathfrak S_{\rm rel}=0
\iff
\forall e:\quad F_e=0\quad\text{and}\quad\sigma_e=\frac12.
}
\]
\end{theorem}

\begin{proof}
Every summand in \(\mathfrak S_{\rm rel}\) is non-negative.  A finite or
countable sum of non-negative terms is zero exactly when every term is zero.
Positive definiteness gives \(\|F_e\|_{G_e}^2=0\iff F_e=0\), while the
information-center theorem gives \(D_H(\sigma_e)=0\iff\sigma_e=1/2\).
\end{proof}

A local \emph{Lagrange node} here means a zero of the declared relational
vector field \(F_e\).  This is an abstract project definition; it does not
assert that every celestial-mechanics Lagrange point is a geometric midpoint.

\section{Third side: reciprocal--conjugation geometry}

Let
\[
s=\sigma+it,
\qquad
\Omega(s)=\frac{s}{1-s}.
\]
The critical line is exactly the projective unit circle:
\[
\Re s=\frac12
\iff
|\Omega(s)|=1.
\]
Define
\[
\Delta_{\rm RC}(u)
=
\left|u^{-1}-\overline u\right|^2.
\]
The formal layer proves, for \(s\ne0,1\),
\[
\boxed{
\Delta_{\rm RC}(\Omega(s))
=
\frac{4(\Re s-1/2)^2}{|s|^2|1-s|^2}.
}
\]
Hence
\[
\boxed{
\Delta_{\rm RC}(\Omega(s))=0
\iff
\Re s=\frac12.
}
\]

Inside the declared binary-spinor model, the phase condition
\[
e^{2\pi i\sigma}=-1
\]
also selects the unique value \(\sigma=1/2\) in the open unit interval.

\begin{corollary}[Three-sided relational zero]
The exchange fixed point, the Shannon/KL zero, and the reciprocal--conjugation
defect zero have the same locus:
\[
\boxed{
x=0
\iff
\sigma=\frac12
\iff
H_2(\sigma)=\ln2
\iff
\Delta_{\rm RC}=0.
}
\]
\end{corollary}

\section{Zero as a relationally forced singular node}

In the project terminology define a relational singular node by
\[
\text{all declared relational defects}=0.
\]
Under the preceding hypotheses this forces
\[
F_e=0,\qquad
D_H(\sigma_e)=0,\qquad
\sigma_e=\frac12
\]
for each participating relation.  Zero is therefore represented as the
centered vanishing of relational imbalance, not as an independently inserted
absolute origin.

\section{Riemann symmetry and the one remaining analytic edge}

The completed zeta symmetry acts horizontally as
\[
\sigma\longmapsto1-\sigma.
\]
Its fixed line is therefore \(\Re s=1/2\).  Symmetry alone does not collapse
an off-axis pair
\[
\frac12+x+it,\qquad\frac12-x+it,
\]
because its midpoint lies on the critical line even when \(x\ne0\).
Consequently an RH proof needs a mechanism forcing the transverse defect
itself to vanish at every non-trivial zero.

\section{Occam minimum proof path}

Let \(E:\C\to\R\) be an independently constructed scalar energy and let
\(c>0\).  Suppose every non-trivial zeta zero \(\rho\) satisfies
\[
E(\rho)=0
\]
and
\[
\boxed{
c\,\Delta_{\rm RC}(\Omega(\rho))
\le
E(\rho).
}
\]
Then
\[
0
\le
c\,\Delta_{\rm RC}(\Omega(\rho))
\le
0,
\]
so
\[
\Delta_{\rm RC}(\Omega(\rho))=0.
\]
The exact projective theorem gives \(\Re\rho=1/2\).  Since \(\rho\) was
arbitrary, RH follows.

This endgame is formalized in
\texttt{SecretOfAHalfFormal/RadialDefect.lean} as
\texttt{riemannHypothesis\_of\_coercive\_zero\_energy}.

The minimum graph is therefore
\[
\boxed{
\xi(\rho)=0
\longrightarrow
E(\rho)=0
\stackrel{E\ge c\Delta}{\longrightarrow}
\Delta(\rho)=0
\longrightarrow
\Re\rho=\frac12.
}
\]

\section{Occam axiomatic RH theorem}

Add to the relational axioms the following single analytic closure axiom.

\begin{axiom}[Canonical zero-energy coercivity]
\label{ax:occam-zero-energy}
There exist a canonical independently defined energy \(E\) and \(c>0\)
such that every non-trivial zeta zero \(\rho\) satisfies
\[
E(\rho)=0,
\qquad
c\,\Delta_{\rm RC}(\Omega(\rho))\le E(\rho).
\]
\end{axiom}

\begin{theorem}[RH in the Occam relational system]
Let \(\mathsf O\) denote the relational axioms together with
\Cref{ax:occam-zero-energy}.  Then
\[
\boxed{\mathsf O\vdash\mathrm{RH}.}
\]
\end{theorem}

\begin{proof}
Apply the coercive zero-energy argument above to an arbitrary non-trivial zero
\(\rho\).  Non-negativity forces the reciprocal--conjugation defect to
vanish, hence \(\Re\rho=1/2\).  Therefore all non-trivial zeros lie on the
critical line.
\end{proof}

Thus the axiomatic proof channel is closed: after the zero-energy coercivity
axiom no further conjectural edge remains.

\section{Proof-cost firewall}

The formal layer also proves
\[
\texttt{ZeroHalfAxisDefectCondition}
\iff
\mathrm{RH}
\]
and
\[
\texttt{ZeroReciprocalDefectCondition}
\iff
\mathrm{RH}.
\]
Therefore simply postulating that the defect vanishes on every non-trivial
zero would be an RH-equivalent assumption.  The Occam theorem above is a
valid proof inside the enlarged axiomatic system, but it becomes an
unconditional proof in ordinary mathematics only when the canonical energy
and its coercive inequality are derived independently of RH or another
RH-equivalent premise.

\section{Minimum proof cost}

The exact crosswalks show that the native half-axis defect, the
reciprocal--conjugation defect, and the Suzuki imaginary-coordinate defect
are three representations of the same transverse quantity.  Therefore this
proof channel does not require simultaneous closure of every historical
route.  Its single independent incoming obligation is:
\[
\boxed{
\text{construct canonical }E
\text{ with }E(\rho)=0
\text{ and }E\ge c\Delta_{\rm RC},
\quad c>0.
}
\]
All subsequent edges are exact and already formalized.

\section{Conclusion}

The relational-zero theorem has three independent faces:
\[
\boxed{
0_{\rm relational}
\equiv
\frac12_{\rm complement}
\equiv
\ln2_{\rm Shannon\ maximum}
\equiv
\Delta_{\rm RC}=0.
}
\]
The half-axis is therefore the unique critical axis of the normalized
relational barycenter.

Within the declared Occam axiomatic system, RH follows with no further gap.
The repository nevertheless retains the external proof firewall: the incoming
canonical zero-energy coercivity theorem must be independently derived before
the result is promoted as an unconditional proof of RH.

\status{The three-sided relational-zero theorem and the coercive zero-energy
RH implication are exact.  The Occam axiomatic channel is closed as an
implication.  Independent derivation of the canonical zero-energy coercivity
edge remains OPEN; therefore the repository does not promote an unconditional
proof of RH.}